\setlength{\absparsep}{18pt} % ajusta o espaçamento dos parágrafos do resumo
\begin{resumo}
O crescimento das pesquisas em finanças quantitativas no ambiente acadêmico tem evidenciado a necessidade de infraestruturas de dados financeiras robustas e de alto desempenho. No âmbito do projeto Análise de Séries Financeiras, do Laboratório de Tecnologias Inovadoras (LTI) da Universidade Federal do Ceará (UFC), Campus de Russas, a falta de uma solução centralizada obrigava pesquisadores a implementar individualmente conexões incovenientes e relativamente complexas com provedores de mercado. Para resolver esse problema, este trabalho apresenta o projeto e o desenvolvimento da plataforma qDATA, uma infraestrutura de dados financeiros de alta velocidade voltada para a aquisição, o armazenamento e a distribuição de dados históricos e em tempo real (barras OHLC). A ingestão de dados ocorre por meio de uma API REST baseada em FastAPI, que armazena os registros em um banco de dados PostgreSQL com a extensão TimescaleDB, otimizada para séries temporais e empregando estratégia de \textit{upsert}. Para o fornecimento de dados aos consumidores finais com baixa latência, o sistema oferece canais gRPC para contextos de alto desempenho, canais HTTPS para protótipos, além de garantir a segurança e controle de acesso granular via JWT e Supabase. O qDATA compõe um ecossistema de serviços financeiros mais amplo, o LTI FinHub. Esse ecossistema baseia-se em uma arquitetura de microsserviços em contêineres Docker e um cliente front-end web. Em vista disso, o qDATA foi pensado a partir de uma arquitetura que separa responsabilidades, portanto tem enfoque total em ingestão, armazenamento e distribuição. O sistema foi validado através de um dashboard analítico por usuários reais, e em se tratando de eliminação de redundâncias o sistema já abriu caminho para que novos pesquisadores do LTI Finanças realizem seus trabalhos sem a dependência do MetaTrader 5. Alguns exemplos incluem, além do (1) dashboard, o serviço de (2) cálculo e simulação de correlações e o serviço de (3) detecção de anomalias em séries temporais.

 \textbf{Palavras-chave}: Séries Temporais Financeiras. Arquitetura de Microsserviços. Docker. TimescaleDB. gRPC. gRPC-Web. HTTPS.
\end{resumo}
